% !TEX spellcheck = en_US



% ~~~~~~~~~~~~~~~~~~~~~~~~~~~~~~~~~~~~~~~ V
% (NC) 2026 Haluk Bingol
% github.com/halukbingol/LaTeX-Templates
%
% Licensees may copy, distribute, display, and perform the work and make derivative works 
% and remixes based on it only for non-commercial purposes. 
% ~~~~~~~~~~~~~~~~~~~~~~~~~~~~~~~~~~~~~~~ A




% =====================================================================
%  Strings in Java -- Learning by Example
%  ---------------------------------------------------------------------
%  Built on hbCisLaTeXTemplate.tex with the libraries in ../../hblib/.
%  Programs and their outputs are read from codeoutput/:
%     codeoutput/<Name>.java   the program
%     codeoutput/<Name>.in     keyboard input (optional)
%     codeoutput/<Name>.txt    what it printed, made by:  sh run_examples.sh
% =====================================================================

\documentclass[aspectratio=169,10pt,compress]{beamer}

% ~~~~~~~~~~~~~~~~~~~~~~~~~~~~~~~~~~~~~~~~~~~~~~~~~~~~~~~~~~~~~~~~ V hblib
	\usepackage{../../hblib/hblibPresentation} 
	\usepackage{../../hblib/hblibCodeoutput}
	\usepackage{../../hblib/hblibDatastructures} 
% ~~~~~~~~~~~~~~~~~~~~~~~~~~~~~~~~~~~~~~~~~~~~~~~~~~~~~~~~~~~~~~~~ A hblib


% ---------------------------------------------------------------------
\title{Strings in Java}
\subtitle{Learning by Example}
\author[Bingol]{Haluk O. Bingol}

\institute[FCIS]{
    Faculty of Computer and Information Sciences\\
    Yeditepe University
}
\date{
	{\tiny \hbTimeStamp}
}

\begin{document}




% ~~~~~~~~~~~~~~~~~~~~~~~~~~~~~~~~~~~~~~ V frm
\begin{frame}[t,fragile]
  \titlepage
\end{frame}
% ~~~~~~~~~~~~~~~~~~~~~~~~~~~~~~~~~~~~~~ A frm




% ~~~~~~~~~~~~~~~~~~~~~~~~~~~~~~~~~~~~~~ V frm
\begin{frame}[t,fragile] \frametitle{Outline}
  \tableofcontents[hideallsubsections]
\end{frame}
% ~~~~~~~~~~~~~~~~~~~~~~~~~~~~~~~~~~~~~~ A frm

% ~~~~~~~~~~~~~~~~~~~~~~~~~~~~~~~~~~~~~~ V
\hSection[Setup]{How to Use These Examples}{}{}
% ~~~~~~~~~~~~~~~~~~~~~~~~~~~~~~~~~~~~~~ A

% ~~~~~~~~~~~~~~~~~~~~~~~~~~~~~~~~~~~~~~ V
\subsection[Run]{Running the examples}
% ~~~~~~~~~~~~~~~~~~~~~~~~~~~~~~~~~~~~~~ A




% ~~~~~~~~~~~~~~~~~~~~~~~~~~~~~~~~~~~~~~ V frm
\begin{frame}[t,fragile] \frametitle{Running the examples}
  \small Every program on these slides is a file in \hCode{codeoutput/}; the output shown is what it really printed (Java 21).
  \begin{columns}[T]
    \begin{column}{0.48\textwidth}
      \begin{block}{macOS (Terminal)}
        \footnotesize
        \hCode{cd} into the folder of the slides, then


% +++++++++++++++++++++++++++++++++++++++ V lst
%: -Lst.
\begin{lstlisting}[style=codesnippet, language=bash]
javac -d build codeoutput/Case.java
java -cp build Case
sh run_examples.sh        # all of them
\end{lstlisting}
% +++++++++++++++++++++++++++++++++++++++ A lst



      \end{block}
    \end{column}
    \begin{column}{0.48\textwidth}
      \begin{block}{Windows (Command Prompt)}
        \footnotesize
        \hCode{chcp 65001} first, so Turkish letters print correctly:


% +++++++++++++++++++++++++++++++++++++++ V lst
%: -Lst.
\begin{lstlisting}[style=codesnippet, language=bash]
chcp 65001
javac -encoding UTF-8 -d build codeoutput\Case.java
java -cp build Case
\end{lstlisting}
% +++++++++++++++++++++++++++++++++++++++ A lst



        \hCode{run\_examples.sh} runs in \emph{Git Bash}.
      \end{block}
    \end{column}
  \end{columns}
  \vspace{0.3em}
  \small Each slide shows the code first; the output appears on the next click. \hRemark{Predict it before you click.}
\end{frame}
% ~~~~~~~~~~~~~~~~~~~~~~~~~~~~~~~~~~~~~~ A frm




% ~~~~~~~~~~~~~~~~~~~~~~~~~~~~~~~~~~~~~~ V
\hSection[Inspect]{Creating and Inspecting}{}{}
% ~~~~~~~~~~~~~~~~~~~~~~~~~~~~~~~~~~~~~~ A

% ~~~~~~~~~~~~~~~~~~~~~~~~~~~~~~~~~~~~~~ V
\subsection[Create]{Creating strings}
% ~~~~~~~~~~~~~~~~~~~~~~~~~~~~~~~~~~~~~~ A




% ~~~~~~~~~~~~~~~~~~~~~~~~~~~~~~~~~~~~~~ V frm
\begin{frame}[t,fragile] \frametitle{Five ways to get a String}
  \lstcoderightstep[0.7][basicstyle=\ttfamily\scriptsize][basicstyle=\ttfamily\tiny]{codeoutput/Create.java}{Java}{codeoutput/Create.txt}
\end{frame}
% ~~~~~~~~~~~~~~~~~~~~~~~~~~~~~~~~~~~~~~ A frm

% ~~~~~~~~~~~~~~~~~~~~~~~~~~~~~~~~~~~~~~ V
\subsection[Length]{Length and characters}
% ~~~~~~~~~~~~~~~~~~~~~~~~~~~~~~~~~~~~~~ A




% ~~~~~~~~~~~~~~~~~~~~~~~~~~~~~~~~~~~~~~ V frm
\begin{frame}[t,fragile] \frametitle{\hCode{length()} and \hCode{charAt(i)}}
  \lstcoderightstep[0.7][basicstyle=\ttfamily\scriptsize][basicstyle=\ttfamily\tiny]{codeoutput/LengthCharAt.java}{Java}{codeoutput/LengthCharAt.txt}

  \vspace{0.3em}
  \small Valid indices are \hCode{0} to \hCode{length()-1}. A \hCode{char} is a number: \hCode{(int) c} shows its code.
\end{frame}
% ~~~~~~~~~~~~~~~~~~~~~~~~~~~~~~~~~~~~~~ A frm

% ~~~~~~~~~~~~~~~~~~~~~~~~~~~~~~~~~~~~~~ V
\subsection[Substring]{Pieces of a string}
% ~~~~~~~~~~~~~~~~~~~~~~~~~~~~~~~~~~~~~~ A




% ~~~~~~~~~~~~~~~~~~~~~~~~~~~~~~~~~~~~~~ V frm
\begin{frame}[t,fragile] \frametitle{\hCode{substring(begin, end)}}
  \begin{columns}[T]
    \begin{column}{0.66\textwidth}
      \lstcode[basicstyle=\ttfamily\scriptsize]{codeoutput/Substring.java}{Java}
    \end{column}
    \begin{column}{0.32\textwidth}
      \begin{center}
        \begin{tikzpicture}\hString[0.32]{H,e,l,l,o,{,},\textvisiblespace,W,o,r,l,d}[7][12]\end{tikzpicture}
      \end{center}
      \small \hCode{substring(7, 12)}: from index 7 \textbf{up to, not including} 12.
      \vspace{0.3em}

      \uncover<2->{\lstoutput[basicstyle=\ttfamily\scriptsize]{codeoutput/Substring.txt}}
    \end{column}
  \end{columns}
\end{frame}
% ~~~~~~~~~~~~~~~~~~~~~~~~~~~~~~~~~~~~~~ A frm

% ~~~~~~~~~~~~~~~~~~~~~~~~~~~~~~~~~~~~~~ V
\hSection[Search]{Searching and Comparing}{}{}
% ~~~~~~~~~~~~~~~~~~~~~~~~~~~~~~~~~~~~~~ A

% ~~~~~~~~~~~~~~~~~~~~~~~~~~~~~~~~~~~~~~ V
\subsection[Search]{Searching}
% ~~~~~~~~~~~~~~~~~~~~~~~~~~~~~~~~~~~~~~ A




% ~~~~~~~~~~~~~~~~~~~~~~~~~~~~~~~~~~~~~~ V frm
\begin{frame}[t,fragile] \frametitle{Searching: \hCode{indexOf}, \hCode{contains}, \dots}
  \lstcoderightstep[0.7][basicstyle=\ttfamily\scriptsize][basicstyle=\ttfamily\scriptsize]{codeoutput/Search.java}{Java}{codeoutput/Search.txt}

  \vspace{0.3em}
  \small \hCode{indexOf} returns \hCode{-1} when not found. The loop finds every occurrence by starting after the previous one.
\end{frame}
% ~~~~~~~~~~~~~~~~~~~~~~~~~~~~~~~~~~~~~~ A frm

% ~~~~~~~~~~~~~~~~~~~~~~~~~~~~~~~~~~~~~~ V
\subsection[Compare]{Comparing}
% ~~~~~~~~~~~~~~~~~~~~~~~~~~~~~~~~~~~~~~ A




% ~~~~~~~~~~~~~~~~~~~~~~~~~~~~~~~~~~~~~~ V frm
\begin{frame}[t,fragile] \frametitle{Comparing: \hCode{equals} and \hCode{compareTo}}
  \lstcoderightstep[0.7][basicstyle=\ttfamily\scriptsize][basicstyle=\ttfamily\scriptsize]{codeoutput/Compare.java}{Java}{codeoutput/Compare.txt}

  \vspace{0.3em}
  \small \hCode{==} asks ``same object?''; \hCode{equals} asks ``same text?''. \hCode{compareTo} is negative, zero or positive.
\end{frame}
% ~~~~~~~~~~~~~~~~~~~~~~~~~~~~~~~~~~~~~~ A frm

% ~~~~~~~~~~~~~~~~~~~~~~~~~~~~~~~~~~~~~~ V
\subsection[Case]{Upper and lower case}
% ~~~~~~~~~~~~~~~~~~~~~~~~~~~~~~~~~~~~~~ A




% ~~~~~~~~~~~~~~~~~~~~~~~~~~~~~~~~~~~~~~ V frm
\begin{frame}[t,fragile] \frametitle{Upper and lower case: the Turkish \hCode{i}}
  \lstcoderightstep[0.7][basicstyle=\ttfamily\scriptsize][basicstyle=\ttfamily\scriptsize]{codeoutput/Case.java}{Java}{codeoutput/Case.txt}

  \vspace{0.3em}
  \small \hRemark{Pitfall:} \hCode{toUpperCase()} without a locale uses the computer's language. On a Turkish
  computer \hCode{"title".toUpperCase()} is \hCode{"TİTLE"}. Use \hCode{Locale.ROOT} for program text (keywords, file names).
\end{frame}
% ~~~~~~~~~~~~~~~~~~~~~~~~~~~~~~~~~~~~~~ A frm

% ~~~~~~~~~~~~~~~~~~~~~~~~~~~~~~~~~~~~~~ V
\hSection[Transform]{Transforming}{}{}
% ~~~~~~~~~~~~~~~~~~~~~~~~~~~~~~~~~~~~~~ A

% ~~~~~~~~~~~~~~~~~~~~~~~~~~~~~~~~~~~~~~ V
\subsection[Split]{Trimming, splitting, joining}
% ~~~~~~~~~~~~~~~~~~~~~~~~~~~~~~~~~~~~~~ A




% ~~~~~~~~~~~~~~~~~~~~~~~~~~~~~~~~~~~~~~ V frm
\begin{frame}[t,fragile] \frametitle{Trimming, splitting and joining}
  \lstcoderightstep[0.7][basicstyle=\ttfamily\scriptsize][basicstyle=\ttfamily\tiny]{codeoutput/TrimSplit.java}{Java}{codeoutput/TrimSplit.txt}

  \vspace{0.3em}
  \small \hCode{split} takes a \textbf{regular expression}: \hCode{"\textbackslash\textbackslash s+"} means ``one or more white-space characters.
\end{frame}
% ~~~~~~~~~~~~~~~~~~~~~~~~~~~~~~~~~~~~~~ A frm

% ~~~~~~~~~~~~~~~~~~~~~~~~~~~~~~~~~~~~~~ V
\subsection[Replace]{Replacing}
% ~~~~~~~~~~~~~~~~~~~~~~~~~~~~~~~~~~~~~~ A




% ~~~~~~~~~~~~~~~~~~~~~~~~~~~~~~~~~~~~~~ V frm
\begin{frame}[t,fragile] \frametitle{Replacing and repeating}
  \lstcoderightstep[0.7][basicstyle=\ttfamily\scriptsize][basicstyle=\ttfamily\scriptsize]{codeoutput/Replace.java}{Java}{codeoutput/Replace.txt}

  \vspace{0.3em}
  \small The last line proves it again: none of these methods changes \hCode{s}; each returns a new string.
\end{frame}
% ~~~~~~~~~~~~~~~~~~~~~~~~~~~~~~~~~~~~~~ A frm

% ~~~~~~~~~~~~~~~~~~~~~~~~~~~~~~~~~~~~~~ V
\subsection[Format]{Formatting}
% ~~~~~~~~~~~~~~~~~~~~~~~~~~~~~~~~~~~~~~ A




% ~~~~~~~~~~~~~~~~~~~~~~~~~~~~~~~~~~~~~~ V frm
\begin{frame}[t,fragile] \frametitle{Formatting numbers and text}
  \lstcoderightstep[0.7][basicstyle=\ttfamily\scriptsize][basicstyle=\ttfamily\scriptsize]{codeoutput/Format.java}{Java}{codeoutput/Format.txt}

  \vspace{0.3em}
  \small \hCode{\%s} text, \hCode{\%d} integer, \hCode{\%.2f} two decimals, \hCode{\%-8s} left-aligned in 8 columns, \hCode{\%n} new line.
  With the Turkish locale the decimal separator is a comma.
\end{frame}
% ~~~~~~~~~~~~~~~~~~~~~~~~~~~~~~~~~~~~~~ A frm

% ~~~~~~~~~~~~~~~~~~~~~~~~~~~~~~~~~~~~~~ V
\subsection[Convert]{Converting}
% ~~~~~~~~~~~~~~~~~~~~~~~~~~~~~~~~~~~~~~ A




% ~~~~~~~~~~~~~~~~~~~~~~~~~~~~~~~~~~~~~~ V frm
\begin{frame}[t,fragile] \frametitle{Strings and numbers}
  \lstcoderightstep[0.7][basicstyle=\ttfamily\scriptsize][basicstyle=\ttfamily\scriptsize]{codeoutput/Convert.java}{Java}{codeoutput/Convert.txt}

  \vspace{0.3em}
  \small \hCode{"42" + 1} is \hCode{"421"}: as soon as one side of \hCode{+} is a \hCode{String}, \hCode{+} means concatenation.
\end{frame}
% ~~~~~~~~~~~~~~~~~~~~~~~~~~~~~~~~~~~~~~ A frm

% ~~~~~~~~~~~~~~~~~~~~~~~~~~~~~~~~~~~~~~ V
\hSection[Build]{Building Strings}{}{}
% ~~~~~~~~~~~~~~~~~~~~~~~~~~~~~~~~~~~~~~ A

% ~~~~~~~~~~~~~~~~~~~~~~~~~~~~~~~~~~~~~~ V
\subsection[Builder]{StringBuilder}
% ~~~~~~~~~~~~~~~~~~~~~~~~~~~~~~~~~~~~~~ A




% ~~~~~~~~~~~~~~~~~~~~~~~~~~~~~~~~~~~~~~ V frm
\begin{frame}[t,fragile] \frametitle{\hCode{StringBuilder}: a mutable string}
  \lstcoderightstep[0.7][basicstyle=\ttfamily\scriptsize][basicstyle=\ttfamily\scriptsize]{codeoutput/Builder.java}{Java}{codeoutput/Builder.txt}

  \vspace{0.3em}
  \small Methods return the builder itself, so calls can be chained: \hCode{sb.append(a).append(b)}.
\end{frame}
% ~~~~~~~~~~~~~~~~~~~~~~~~~~~~~~~~~~~~~~ A frm

% ~~~~~~~~~~~~~~~~~~~~~~~~~~~~~~~~~~~~~~ V
\subsection[Words]{Working word by word}
% ~~~~~~~~~~~~~~~~~~~~~~~~~~~~~~~~~~~~~~ A




% ~~~~~~~~~~~~~~~~~~~~~~~~~~~~~~~~~~~~~~ V frm
\begin{frame}[t,fragile] \frametitle{Word by word: title case}
  \lstcoderightstep[0.7][basicstyle=\ttfamily\scriptsize][basicstyle=\ttfamily\scriptsize]{codeoutput/TitleCase.java}{Java}{codeoutput/TitleCase.txt}
\end{frame}
% ~~~~~~~~~~~~~~~~~~~~~~~~~~~~~~~~~~~~~~ A frm

% ~~~~~~~~~~~~~~~~~~~~~~~~~~~~~~~~~~~~~~ V
\subsection[Chars]{Classifying characters}
% ~~~~~~~~~~~~~~~~~~~~~~~~~~~~~~~~~~~~~~ A




% ~~~~~~~~~~~~~~~~~~~~~~~~~~~~~~~~~~~~~~ V frm
\begin{frame}[t,fragile] \frametitle{Classifying characters: \hCode{Character}}
  \lstcoderightstep[0.7][basicstyle=\ttfamily\scriptsize][basicstyle=\ttfamily\scriptsize]{codeoutput/CharClass.java}{Java}{codeoutput/CharClass.txt}

  \vspace{0.3em}
  \small \hCode{Character.isLetter} knows Turkish letters too. \hCode{\textquotesingle 7\textquotesingle{} - \textquotesingle 0\textquotesingle} turns a digit character into its value.
\end{frame}
% ~~~~~~~~~~~~~~~~~~~~~~~~~~~~~~~~~~~~~~ A frm

% ~~~~~~~~~~~~~~~~~~~~~~~~~~~~~~~~~~~~~~ V
\hSection[Modern]{Modern Java}{}{}
% ~~~~~~~~~~~~~~~~~~~~~~~~~~~~~~~~~~~~~~ A

% ~~~~~~~~~~~~~~~~~~~~~~~~~~~~~~~~~~~~~~ V
\subsection[Text block]{Text blocks}
% ~~~~~~~~~~~~~~~~~~~~~~~~~~~~~~~~~~~~~~ A




% ~~~~~~~~~~~~~~~~~~~~~~~~~~~~~~~~~~~~~~ V frm
\begin{frame}[t,fragile] \frametitle{Text blocks (Java 15+)}
  \lstcoderightstep[0.7][basicstyle=\ttfamily\scriptsize][basicstyle=\ttfamily\scriptsize]{codeoutput/TextBlock.java}{Java}{codeoutput/TextBlock.txt}

  \vspace{0.3em}
  \small A text block starts with \hCode{"""} and a new line. The common indentation is removed;
  a \hCode{\textbackslash} at the end of a line joins it with the next one.
\end{frame}
% ~~~~~~~~~~~~~~~~~~~~~~~~~~~~~~~~~~~~~~ A frm

% ~~~~~~~~~~~~~~~~~~~~~~~~~~~~~~~~~~~~~~ V
\subsection[Switch]{Switch on strings}
% ~~~~~~~~~~~~~~~~~~~~~~~~~~~~~~~~~~~~~~ A




% ~~~~~~~~~~~~~~~~~~~~~~~~~~~~~~~~~~~~~~ V frm
\begin{frame}[t,fragile] \frametitle{\hCode{switch} on strings}
  \lstcoderightstep[0.7][basicstyle=\ttfamily\scriptsize][basicstyle=\ttfamily\scriptsize]{codeoutput/SwitchString.java}{Java}{codeoutput/SwitchString.txt}

  \vspace{0.3em}
  \small A \hCode{switch} can test strings (with \hCode{equals}, never \hCode{==}). Normalize the case first, with \hCode{Locale.ROOT}.
\end{frame}
% ~~~~~~~~~~~~~~~~~~~~~~~~~~~~~~~~~~~~~~ A frm

% ~~~~~~~~~~~~~~~~~~~~~~~~~~~~~~~~~~~~~~ V
\hSection[Summary]{Summary}{}{}
% ~~~~~~~~~~~~~~~~~~~~~~~~~~~~~~~~~~~~~~ A

% ~~~~~~~~~~~~~~~~~~~~~~~~~~~~~~~~~~~~~~ V
\subsection[Cheat sheet]{Cheat sheet}
% ~~~~~~~~~~~~~~~~~~~~~~~~~~~~~~~~~~~~~~ A




% ~~~~~~~~~~~~~~~~~~~~~~~~~~~~~~~~~~~~~~ V frm
\begin{frame}[t,fragile] \frametitle{Cheat sheet}
  \footnotesize
  \begin{columns}[T]
    \begin{column}{0.48\textwidth}
      \begin{tabular}{@{}L{0.5\textwidth}L{0.42\textwidth}@{}}
        \toprule
        \textbf{Method} & \textbf{Gives} \\
        \midrule
        \hCode{length()}                 & number of chars \\
        \hCode{charAt(i)}                & char at index \hCode{i} \\
        \hCode{substring(b, e)}          & chars \hCode{b..e-1} \\
        \hCode{indexOf(x)}               & first index or \hCode{-1} \\
        \hCode{contains(x)}              & \hCode{boolean} \\
        \hCode{startsWith}/\hCode{endsWith} & \hCode{boolean} \\
        \hCode{equals}/\hCode{equalsIgnoreCase} & same text? \\
        \hCode{compareTo(x)}             & $<0$, $0$, $>0$ \\
        \bottomrule
      \end{tabular}
    \end{column}
    \begin{column}{0.48\textwidth}
      \begin{tabular}{@{}L{0.5\textwidth}L{0.42\textwidth}@{}}
        \toprule
        \textbf{Method} & \textbf{Gives (a new string)} \\
        \midrule
        \hCode{toUpperCase(locale)}      & upper case \\
        \hCode{strip()}                  & no outer spaces \\
        \hCode{split(regex)}             & \hCode{String[]} \\
        \hCode{String.join(sep, ...)}    & joined string \\
        \hCode{replace(a, b)}            & all \hCode{a} replaced \\
        \hCode{repeat(n)}                & \hCode{n} copies \\
        \hCode{String.format(...)}       & formatted text \\
        \hCode{Integer.parseInt(s)}      & \hCode{int} (not a string) \\
        \bottomrule
      \end{tabular}
    \end{column}
  \end{columns}
\end{frame}
% ~~~~~~~~~~~~~~~~~~~~~~~~~~~~~~~~~~~~~~ A frm

% ~~~~~~~~~~~~~~~~~~~~~~~~~~~~~~~~~~~~~~ V
\subsection[Exercises]{Exercises}
% ~~~~~~~~~~~~~~~~~~~~~~~~~~~~~~~~~~~~~~ A




% ~~~~~~~~~~~~~~~~~~~~~~~~~~~~~~~~~~~~~~ V frm
\begin{frame}[t,fragile] \frametitle{Exercises}
  \begin{enumerate}\small
    \item Print a name in reverse with a loop, then with \hCode{StringBuilder.reverse()}.
    \item Count the vowels (\hCode{a e ı i o ö u ü}) in a Turkish sentence.
    \item Given \hCode{"surname, name"}, print \hCode{"Name SURNAME"} (use \hCode{indexOf}, \hCode{substring}, \hCode{strip}).
    \item Read a CSV line \hCode{"Ali,3.25,6"} and print it as a formatted table row with \hCode{String.format}.
    \item Check whether a word is a palindrome, ignoring case (\hCode{"Kayak"}).
    \item Turn \hCode{"hello world from java"} into \hCode{"helloWorldFromJava"} (camel case).
    \item Explain the output of \hCode{System.out.println(1 + 2 + "3" + 4 + 5);} and then run it.
  \end{enumerate}
\end{frame}
% ~~~~~~~~~~~~~~~~~~~~~~~~~~~~~~~~~~~~~~ A frm




% ~~~~~~~~~~~~~~~~~~~~~~~~~~~~~~~~~~~~~~ V frm
\begin{frame}[t,fragile]
	\begin{center}
		\vspace{4cm}
		\hRemark{\Huge Questions?}
	\end{center}
\end{frame}
% ~~~~~~~~~~~~~~~~~~~~~~~~~~~~~~~~~~~~~~ A frm



\end{document}
